\section{Topics, Organization, and Timeline}

We will invite paper submissions in the following four topic areas to meet the workshop goals. However, we will encourage submission on uncertainty visualization that may not necessarily fit into the following topic areas. 

\vspace{-6pt}
\begin{itemize}[noitemsep]
    \item 
    \textbf{Uncertainty Propagation}: Techniques to effectively model data uncertainty, Statistical/Probabilistic techniques to propagate uncertainty through complex data processing pipelines, Uncertainty propagation techniques in application domain, including but not limited to climate science, quantum computing, energy science, and material science.
    
    \item 
    \textbf{Uncertainty Communication}: Techniques to convey uncertainty in 2D/3D/high-dimensional (scientific or information) data, Effective ways to educate non-experts and scientific users about uncertainty visualization, Visualization tools or software to communicate uncertainty. 

    \item 
    \textbf{AI and Uncertainty Visualization}: Techniques to visualize uncertainty in AI models for explainability, AI techniques for propagation and visualization of uncertainty in data, and AI for generalization of uncertainty visualization across application domains.

     \item 
     \textbf{Decision-Making Under Uncertainty}: Obstacles to effective decision-making under uncertainty, User studies to evaluate the influence of uncertainty visualization on decision-making, and AI and statistical techniques for robust decision-making under uncertainty.
\end{itemize}

\vspace{-4pt}
{\large Paper Types:} $\;$ We will accept short papers with 4 pages (plus 1 page of references) and long papers with 9 pages (plus 1 page of references) on the aforementioned topics or other topics related to uncertainty. We are interested in including these papers in {\em regular publication category} listed on the VIS 2025 call for workshops.
%

{\large Deadlines.} $\;$ The submission deadline for papers will be several weeks after the VIS author notification date (i.e., June 06, 2025). The tentative deadlines are as follows.

\vspace{2pt}
\indent\indent\textbf{Call for papers and website in action:} April 30, 2025\\
\indent\indent\textbf{Deadline for submission:} July 10, 2025 \\
\indent\indent\textbf{Notification of acceptance:} July 31, 2025 \\
\indent\indent\textbf{Camera-ready papers due:} August 18, 2025 
%

\vspace{2pt}
{\large Process.} $\;$ We will set up a small program committee of $15$-$20$ members (similar to last year) to have 2-3 reviewers on each submitted paper. If the number of submissions is lower than anticipated, we will then distribute the reviewing task among the workshop organizers. The final determination on which papers are accepted will be done based on individual review scores and thorough discussion among the workshop organizers. The workshop organizers will select from the accepted submissions a list of best paper candidates (typically those with the highest review scores). Again, the best paper and any honorable mentions will be selected based on the through discussion among the workshop organizers.
%

{\large Responsibilities.} $\;$ Dr. Athawale will be the point of contact and overall chair for this workshop, and others will be on the workshop committee. Our main tasks will include coordinating with VIS, organizing a panel, setting a budget, distributing paper calls, announcing the program, managing the workshop website, and recruiting sponsors. The workshop organizers will be responsible for the timely delivery of camera-ready revisions of accepted papers to the VIS publication chair. We will finalize the panel members (and keynote speakers in the event of execution of the back-up plan [see Sec.~\ref{sec:backup}]) by August 18, 2025. 
%The details regarding the selection of panel members is in Sec.~\ref{sec:activities}.\\

{\large Diversity in Participation} $\;$ We will invite members of the panel session (the keynote/external speakers in the case of back-up plan) to ensure diversity in terms of their race, gender, and background (similar to last year's workshop). The workshop organizers will reach out to non-R1 universities, minority serving institutions (MSI), and diversity-focused organizations (National Society of Black Engineers, American Association of People with Disabilities, WomenInHPC) to encourage submissions of papers.  


% We will invite submissions in the following topic areas for papers and posters. The topic areas touch upon the applications, techniques, software, and decision-making aspects of uncertainty visualization to meet the workshop goals.

% \begin{itemize}
%     \item Applications: Domains-specific use cases for uncertainty visualization and analysis of scientific or information data with domains including but not limited to climate science, energy science, material science, and machine learning.

%     \item Techniques: Techniques to quantify and convey uncertainty in 2D/3D/high-dimensional (scientific or information) data. These cover a broad range of concepts, including but not limited to statistics, machine learning,  computing, rendering, perception, and cognition that are fundamental to uncertainty analysis.

%     \item Software and tools: Uncertainty visualization software or tools  to convey the uncertainty in general or application-specific scientific or information data.

%      \item Decision frameworks: Methods, workflows, and application use cases for robust decision-making under uncertainty
% \end{itemize}

% %\ta{Invited domain expert talks/keynote speaker:} \\
% %A speaker who can possibly talk about uncertainty visualization challenges in their application domain.\\\\
% {\large Paper Types.} $\;$ We will accept short papers with 4 pages (plus 1 page of references) that cover a broad range of studies in uncertainty analysis, including novel research contributions, domain-specific or general requirements for successful uncertainty analysis, obstacles to understanding data uncertainty presented through use cases, and successful uncertainty visualization workflows for robust solutions. We are interested to include these papers in {\em regular publication category} listed on the VIS 2024 workshop website.\\ 

% {\large Deadlines.} $\;$ The submission deadline for papers will be several weeks after the VIS author notification date. The tentative deadlines are as follows.\\
% \textbf{ $\;$ Call for papers and website in action:} April 30, 2025\\
% \textbf{Deadline for submission:} July 10, 2025 \\
% \textbf{Notification of acceptance:} July 31, 2025 \\
% \textbf{Camera-ready papers due:} August 18, 2025 \\

% {\large Process.} $\;$ Given a growing interest in uncertainty visualization area and this being the third VIS workshop on uncertainty, we anticipate a reasonable number of submissions. We will therefore set up a small program committee of 5-10 members to review the papers. If the number of submissions is lower than anticipated, we will then distribute the reviewing task among the workshop organizers. The final determination on which papers are accepted will be done based on individual review scores and thorough discussion among the workshop organizers. The workshop organizers will select from the accepted submissions a list of best paper and poster candidates (typically those with the highest review scores). Again, the best paper and any honorable mentions will be selected based on the through discussion among the workshop organizers.\\

% {\large Responsibilities.} $\;$ Dr. Athawale will be the point of contact and overall chair for this workshop, and others will be on the workshop committee. Our main tasks will include coordinating with VIS, identifying the keynote speaker, organizing a panel, setting a budget, distributing paper calls, announcing the program, managing the workshop website, and recruiting sponsors. The workshop organizers will be responsible for the timely delivery of camera-ready revisions of accepted papers to the VIS publication chair. We will finalize the keynote speakers and panel members by August 18, 2024. The details regarding the selection of keynote speakers and panel members is in Sec.~\ref{sec:activities}.\\

% {\large Diversity in Participation} $\;$ We will invite the keynote/external speakers and members of the panel session to ensure diversity in terms of their race, gender, and background. The workshop organizers will reach out to non-R1 universities, minority serving institutions (MSI), and diversity-focused organizations (National Society of Black Engineers, American Association of People with Disabilities, WomenInHPC) to encourage submissions of papers/posters.  

%and provide the manuscripts and all reviews in an anonymized form to an ad-hoc best paper committee (recruited by the paper chairs among the PC) to select the best paper and any honorable mentions. 


%\section {Organizing Committee Roles and Leadership Structure}
%The following are the core roles of the organizing committee:
%\begin{itemize}
%\item Workshop Chairs: The three workshop chairs coordinate with VIS, identify the keynote speaker, organize a panel, set a budget, manage the workshop website, and recruit sponsors.
%\item Paper/Program Chairs: The two paper chairs recruit the program committee, write the call for papers, set up the reviewing site, assign reviewers, determine the accepted papers based on reviews, notify authors of the decisions, recruit a best paper committee, determine the program schedule, and invite session chairs for the onsite program. Furthermore, paper chairs are responsible for the timely delivery of camera-ready revisions of accepted papers to the VIS publication chair.
%\item Communications Chair: The communications chair is responsible for all community communications, including distributing paper calls, announcing the program, and maintaining the workshop website.
%\end{itemize}

%Below are our proposed assignment of roles (tentative):
%\begin{itemize}
%\item	Workshop Chairs: Tushar M. Athawale, Chris R. Johnson, Kristi Potter
%\item	Paper Chairs: Paul Rosen, Dave Pugmire
%\item	Communications Chair: Tushar M. Athawale
%\end{itemize}